\documentclass{article}

\usepackage{amsmath}
\usepackage{amsfonts}
\usepackage{amsthm}
\newtheorem{definition}{Definition}
\newtheorem{theorem}{Theorem}
\newtheorem{lemma}{Lemma}
\newtheorem{claim}{Claim}
\usepackage{authblk}
\usepackage{cryptocode}
\usepackage{csquotes}
\usepackage{enumerate}
\usepackage{float}
\usepackage{framed}
\usepackage[colorlinks]{hyperref}
\def\tmp#1#2#3{
  \definecolor{Hy#1color}{#2}{#3}
  \hypersetup{#1color=Hy#1color}}
\tmp{link}{HTML}{800006}
\tmp{cite}{HTML}{2E7E2A} % better default colors for hyperref links; see https://tex.stackexchange.com/a/525297
\usepackage[capitalise]{cleveref}
\usepackage{soul} % allow line break in underline

\newcommand{\group}{\mathbb{G}}
\newcommand{\Z}{\mathbb{Z}}

\newcommand{\negl}{\mathsf{negl}}

\newcommand{\adv}{\mathcal{A}}
\newcommand{\env}{\mathcal{Z}}
\renewcommand{\sim}{\mathcal{S}}
\newcommand{\red}{\mathcal{R}}
\newcommand{\func}{\mathcal{F}}
\newcommand{\Fpake}{\func_\mathrm{PAKE}}

\newcommand{\Adv}{\mathbf{Adv}}
\newcommand{\Dist}{\mathbf{Dist}}

\newcommand{\Gen}{\mathsf{Gen}}
\newcommand{\Encap}{\mathsf{Encap}}
\newcommand{\Decap}{\mathsf{Decap}}

\newcommand{\PC}{\mathsf{PC}}

\newcommand{\Hkey}{H_\mathsf{key}}
\newcommand{\Htag}{H_\mathsf{tag}}

\newcommand{\pw}{\mathsf{pw}}

\newcommand{\NewSession}{\mathsf{NewSession}}
\newcommand{\TestPwd}{\mathsf{TestPwd}}
\newcommand{\NewKey}{\mathsf{NewKey}}

\newcommand{\fresh}{\mathsf{fresh}}
\newcommand{\compromised}{\mathsf{compromised}}
\newcommand{\interrupted}{\mathsf{interrupted}}
\newcommand{\completed}{\mathsf{completed}}

\newcommand{\qkey}{q_\mathsf{key}}
\newcommand{\qtag}{q_\mathsf{tag}}

\newcommand{\Query}{\mathsf{Query}}

\interfootnotelinepenalty=10000 % force footnotes to appear on the same page

\begin{document}

\title{The Most Efficient Protocol for PAKE: \\
\medskip
\large{What Exact Stuff Do You Need to Hash at the End?}}

\author{Jiayu Xu}
\affil{Oregon State University, \texttt{xujiay@oregonstate.edu}}
\date{}

\maketitle

\section{Introduction}

\section{Preliminaries}
\label{sec:prelim}

\section{Security Analysis of the ``Hash Everything'' Version}

\subsection{The Simulator}
\subsubsection{Simulation Strategy}
A large portion of the simulator is routine; below we explain the idea behind the non-trivial parts.
\paragraph{Simulating honest $P$-to-$P'$ message.}
When $P$'s session begins, $\sim$ must simulate its message $(s,T)$ to $P'$. This can be easily done: just sample $(s,T)$ at random. However, every time $\adv$ queries $H_2(\pw^*,T)$ (let the result be $t^*$), $\sim$ needs to generate a KEM key pair $(pk^*,sk^*)$ for later use (see the ``Extracting password guess from malicious $P'$'' paragraph below) and make sure that $(pk^*,sk^*)$ is consistent with the RO queries. This can be done as follows: $\sim$ computes
\[
r^* := s \oplus t^*, \quad R^* := T/pk^*,
\]
and programs $H_1(\pw^*,r^*) := R^*$. This also allows $\sim$ to store the association between $\pw^*$ and $(pk^*,sk^*)$. (Note that $(pk^*,sk^*)$ needs to be freshly generated every time $\adv$ queries $H_2(\star,T)$; otherwise $R^*$ never changes and $\adv$ can easily find collisions in $H_1$.)
\paragraph{Extracting password guess from malicious $P$.}
When the adversary $\adv$ sends a message $(s^*,T^*)$ to $P'$, if it is different from the honest message $(s,T)$, the simulator $\sim$ must extract a password guess (if any) and send a corresponding $\TestPwd$ command to $\Fpake$. There are two types of attack, which need to be addressed separately:
\begin{enumerate}
  \item \emph{``Normal'' online attack:} $\adv$ executes $P$'s algorithm on a password guess $\pw^*$. That is, $\adv$ samples public key $pk^*$ and randomness $r^*$, computes
      \begin{align*}
      R^* &:= H_1(\pw^*,r^*), &T^* &:= pk^* \cdot R^*, \\
      t^* &:= H_2(\pw^*,T^*), &s^* &:= t^* \oplus r^*,
      \end{align*}
      and sends $(s^*,T^*)$ to $P'$. \\
      It takes a while to realize how to extract $\pw^*$ from $(s^*,T^*)$. First, $\sim$ should scan all $H_2(\star,T^*)$ queries, which yields multiple $\pw^*$ values (and multiple $t^*$ values); $\sim$ must decide which one is the ``actual'' password guess. This can be done via the following. For each such $(\pw^*,t^*)$ pair, $\sim$ computes the corresponding $r^* := s^* \oplus t^*$, and checks if (and when) $H_1(\pw^*,r^*)$ was queried. \emph{With overwhelming probability, there is at most one such $H_1$ query made before the $H_2(\pw^*,T^*)$ query}, from which $\sim$ can extract the ``actual'' password guess $\pw^*$. (This special $H_1(\pw^*,r^*)$ query is called the ``anchor query'' in \cite{CCS:McQRosRoy20,EC:JanRoyXu25}.)
  \item \emph{Related key attack:} This attack works only after $\adv$ receive an honest $(s,T)$ pair from $P$. $\adv$ skips sampling $pk^*$ or $r^*$; instead, $\adv$ picks any $\Delta \in \group$, computes $T^* := T \cdot \Delta$ and
      \begin{align*}
      t^* &:= H_2(\pw^*,T), &(t^*)' &:= H_2(\pw^*,T^*), \\
      \delta &:= t^* \oplus (t')^*, &s^* &:= s \oplus \delta,
      \end{align*}
      and sends $(s^*,T^*)$ to $P'$. \\
      Let us first explain what $\adv$ is trying to achieve here. The same randomness $r$ yields two valid messages, $(s,T)$ and $(s^*,T^*)$, such that
      \[
      s \oplus s^* = H_2(\pw,T) \oplus H_2(\pw,T^*),
      \]
      which allows $\adv$ to pick $T^*$ and ``reverse engineer'' the corresponding $s^*$ if it guesses $\pw$ correctly. $(s^*,T^*)$ implicitly encrypts $pk^* = T^*/H_1(\pw,r)$, which $\adv$ can compute as $pk \cdot (T^*/T)$. \\
      To extract $\pw^*$ from $(s^*,T^*)$, $\sim$ should scan all $H_2(\star,T)$ and $H_2(\star,T^*)$ queries (which yield multiple candidate $\pw^*$ values) and check which ones satisfy
      \[
      H_2(\pw^*,T) \oplus H_2(\pw^*,T^*) = s \oplus s^*.
      \]
      With overwhelming probability, at most one $\pw^*$ will satisfy this equation.
\end{enumerate}
\paragraph{Extracting password guess from malicious $P'$.}
When the adversary $\adv$ sends a message $(c^*,\tau^*)$ to $P$, if $\adv$ is not passive (i.e., it is not the case that $(s^*,T^*) = (s,T)$ \emph{and} $(c^*,\tau^*) = (c,\tau)$), the simulator $\sim$ must extract a password guess (if any) and send a corresponding $\TestPwd$ command to $\Fpake$.

$\sim$ should find the $\Htag$ query whose result is $\tau^*$, which includes $(\pw^*,pk^*,K^*)$. However, $\pw^*$ constitutes a password guess only if it is consistent with $pk^*$ and $K^*$, i.e., $K^* = \Decap(sk^*,c^*)$ for $sk^*$ corresponding to $pk^*$. $\sim$ can check this condition because it has stored the correspondence between $\pw^*$ and $(pk^*,sk^*)$ (see the ``Simulating honest $P$-to-$P'$ message'' paragraph above), and can use $sk^*$ to check if the equation holds. (This extraction strategy does not work anymore if \emph{any} of $\pw^*,pk^*,c^*$ is not included in $\Htag$. In later sections we will explain how to simulate other versions of the protocol.)

\subsubsection{Formal Description of the Simulator}
\label{sec:sim}

\subsection{Hybrid Argument}
We now prove that for any PPT environment $\env$, the simulator $\sim$ in \Cref{sec:sim} generates an ideal-world view indistinguishable from $\env$'s real-world view. The sequence of games start from the real world (Game 0) and ends at the ideal world.

\medskip\noindent\emph{Game 0:} This is the real world. Recall that the passwords of $P$ and $P'$ are $\pw$ and $\pw'$, respectively; the flow of events is:
\begin{itemize}
  \item $P$ sends $(s,T)$ to $P'$;
  \item It is intercepted by the man-in-the-middle adversary $\adv$, who sends $(s^*,T^*)$ (which may or may not be equal to $(s,T)$) to $P'$;
  \item $P'$ outputs session key $SK'$ and sends $(c,\tau')$ to $P$;
  \item It is again intercepted by $\adv$, who sends $(c^*,\tau^*)$ to $P$;
  \item $P$ computes $\tau$. If $\tau^* = \tau$ then $P$ outputs $SK$, otherwise $P$ outputs $\bot$.
\end{itemize}
(Note that there are three $\tau$ values: $\tau'$ computed and sent by $P'$, $\tau^*$ sent by $\adv$, and $\tau$ computed by $P$ used in $P$'s check.)
\paragraph{$\adv$ is passive, except maybe modifying the tag.}
\mbox{}

\medskip\noindent\emph{Game 1:} In the case where $\pw = \pw' \land (s^*,T^*) = (s,T) \land c^* = c$, do the following:
\begin{itemize}
  \item If $\tau^* = \tau'$ (i.e., $\adv$ is passive), set $SK := SK'$.
  \item If $\tau^* \neq \tau'$, set $SK := \bot$.
\end{itemize}
By correctness of the protocol (see \Cref{sec:prelim}), in Game 0 $\tau = \tau'$ except with probability $\Adv^\text{correctness}$. This means that except with probability $\Adv^\text{correctness}$, $P$'s check $\tau^* \stackrel{?}{=} \tau$ in Game 0 is equivalent to $P$'s check $\tau^* \stackrel{?}{=} \tau'$ in Game 1. If the check passes then both games set $SK := \Hkey(\pw,pk,c,K)$, otherwise both games set $SK := \bot$. We have
\[
\Dist_\env^{0,1} = \Adv^\text{correctness}.
\]

\emph{{\color{red}This game corresponds to games $\mathsf{G}_3$--$\mathsf{G}_5$ in \cite{EPRINT:ArrBarJar25}. We believe having a single game that handles both the $\tau^* = \tau'$ case and the $\tau^* \neq \tau'$ case is much cleaner.}}

\paragraph{$P$ and $P'$'s passwords match, and $\adv$ forwards the $P$-to-$P'$ message.}
\mbox{}

\medskip\noindent\emph{Game 2:} In the case where $\tau^* \neq \tau'$, set $SK := \bot$ if either of the following happens:
\begin{itemize}
  \item $\adv$ never makes an $\Htag(\pw,pk,c^*,\star)$ query resulting in $\tau^*$, or
  \item $\adv$ makes two different $\Htag(\pw,pk,c^*,\star)$ queries, both resulting in $\tau^*$.
\end{itemize}

In the first event, Game 1 and Game 2 differ if $\adv$ never makes an $\Htag(\pw,pk,c^*,\star)$ query resulting in $\tau^*$, yet $P$'s check $\tau^* \stackrel{?}{=} \Htag(\pw,pk,c^*,K)$ passes. Clearly,
\begin{itemize}
  \item $\adv$ does not query $\Htag(\pw,pk,c^*,K)$ (otherwise the result is $\tau^*$), and
  \item $P'$ does not query $\Htag(\pw,pk,c^*,K)$ (otherwise $\tau' = \Htag(\pw,pk,c^*,K) = \tau^*$).
\end{itemize}
So $\Htag(\pw,pk,c^*,K)$ is independent of $\adv$'s view, and the probability that $\adv$ sends $\tau^* = \Htag(\pw,pk,c^*,K)$ is $1/2^n$. The second event is collision in $\Htag$, whose probability is upper-bounded by $\qtag(\qtag-1)/2^{n+1}$. Overall,
\[
\Dist_\env^{1,2} \leq \frac{\qtag(\qtag-1)+2}{2^{n+1}}.
\]

\emph{{\color{red}This game corresponds to games $\mathsf{G}_6$ and $\mathsf{G}_7$ in \cite{EPRINT:ArrBarJar25}, except that the $\Htag$ collision case was ruled out in game $\mathsf{G}_2$ in \cite{EPRINT:ArrBarJar25} (as aborting condition [A1]).}}

\medskip\noindent\emph{Game 3:} In the case where $\pw = \pw' \land (s^*,T^*) = (s,T)$, sample  $SK',\tau' \gets \{0,1\}^n$. ($c$ is still computed as before.)

Let $\Query_1$ be the event that $\adv$ queries \emph{either} $\Hkey(\pw,pk,c,K')$ \emph{or} $\Htag(\pw,\allowbreak pk,c,K')$. Clearly Game 2 and Game 3 are identical unless $\Query_1$ happens. We upper-bound $\Pr[\Query_1]$ via a reduction $\red_\text{OW-1PCA1}$ to the OW-1PCA property of KEM:

\medskip\noindent\textbf{\underline{Reduction $\red_\mathrm{OW\text{-}1PCA1}$}:} \quad \textcolor[gray]{0.5}{// \fbox{Boxed text} indicates the single $\PC$ oracle query; same for reductions below}

\medskip\noindent
Sample $i \gets [\qkey+\qtag]$ as a guess that $\adv$'s $i$-th $\Hkey$ or $\Htag$ query triggers $\Query_1$. \\
On input $(pk,c)$,
\begin{enumerate}
  \item Invoke $\env$. On $(\NewSession, sid, P, P', \pw, \text{``initiator''})$ from $\env$, run $P$'s algorithm to compute $(s,T)$, except that $\red_\text{ANO-1PCA}$ uses its input $pk$ rather than $(pk,\star) \gets \Gen$. (Note that computing $(s,T)$ does not require knowing $sk$.) Send $(s,T)$ to $\adv$.
  \item On $(\NewSession, sid, P', P, \pw', \text{``respondent''})$ from $\env$ and $(s^*,T^*)$ from $\adv$, if $\neg(\pw = \pw' \land (s^*,T^*) = (s,T))$, abort. Otherwise sample $SK',\tau' \gets \{0,1\}^n$, output $SK'$ to $\env$, and send $(c,\tau')$ to $\adv$. ($c$ is part of $\red_\text{OW-PCA1}$'s inputs.)
  \item On $(c^*,\tau^*)$ from $\adv$,
      \begin{enumerate}
        \item If $c^* = c$, do what's described in Game 1.
        \item If $c^* \neq c$, find $\adv$'s $\Htag(\pw,pk,c^*,K^*)$ query whose result is $\tau^*$. If there is none or more than one, output $SK := \bot$ to $\env$. (This matches Game 2.)
        \item \fbox{Query $\PC(sk,c^*,K^*)$}. \quad // \textcolor[gray]{0.5}{This is a valid query since we do this step only if $c^* \neq c$} \\
            If the result is 1, output $SK := \Hkey(\pw,pk,c^*,K^*)$ to $\env$. \\
            If the result is 0, output $SK := \bot$ to $\env$.
      \end{enumerate}
  \item On $\adv$'s $i$-th $\Hkey$ or $\Htag$ query, if it is in the form of $\Hkey(\pw,pk,c,(K')^*)$ or $\Htag(\pw,pk,c,(K')^*)$ for some $(K')^*$, output $(K')^*$; otherwise abort.
\end{enumerate}

Assuming $\Query_1$ happens, the only difference between Game 2/3 and $\red_\text{OW-1PCA1}$'s simulation (until $\Query_1$ happens) is that $\red_\text{OW-1PCA1}$ uses its input $pk$ on the $P$ side, and then uses its input $c \gets \Encap(pk)$ on the $P'$ side. Since we assume $\pw = \pw' \land (s^*,T^*) = (s,T)$, in Game 2/3 $P'$'s algorithm will recover $pk' = pk$, so there $c$ is also computed as $\Encap(pk)$ and thus $\red_\text{OW-1PCA1}$ simulates Game 2/3 perfectly. Furthermore, if $\red_\text{OW-1PCA1}$'s guess is correct then it wins the OW-1PCA game. We have
\[
\Dist_\env^{2,3} \leq \Pr[\Query_1] \leq (\qkey+\qtag)\Adv_{\red_\text{OW-1PCA1}}^\text{OW-1PCA}.
\]

\emph{{\color{red}This game corresponds to game $\mathsf{G}_8$ in \cite{EPRINT:ArrBarJar25}. However, there are two differences:
\begin{enumerate}
  \item \cite{EPRINT:ArrBarJar25} reduces to OW-PCA (where the KEM adversary can query its plaintext-checking oracle polynomially many times), rather than OW-1PCA. In this case, step 4 of the reduction above can query $\PC(sk,c,(K')^*)$ $\qkey+\qtag$ times to check which of $\adv$'s query (if any) triggers $\Query_1$, hence avoiding the $\qkey+\qtag$ loss in its advantage.
  \item The above also means that the plaintext-checking oracle in their OW-PCA game has to allow for $\PC(sk,c,\star)$ queries, whereas our reduction to OW-1PCA never makes a query in this form and thus the plaintext-checking oracle in our OW-1PCA game can disallow such queries.
\end{enumerate}
}}

\medskip\noindent\emph{Game 4:} Again in the case where $\pw = \pw' \land (s^*,T^*) = (s,T)$, generate $P'$'s KEM public key $(pk',\star) \gets \Gen$ and compute $(c,\star) \gets \Encap(pk')$. (Note that $K'$ is not used anywhere since Game 3 --- in particular, $P'$ does not use $K'$ to compute $SK'$ or $\tau'$ anymore --- so we do not need to explicitly mention it.)

The difference between Game 3 and Game 4 is that in Game 3 both $P$ and $P'$ use the same $pk$, whereas in Game 4 $P$ uses $pk$ and $P'$ uses an independent $pk'$. We construct a reduction $\red_\text{ANO-1PCA}$ to the ANO-1PCA property of KEM. This reduction, in fact, is identical to the reduction $\red_\text{OW-1PCA1}$ to the OW-1PCA property in Game 3 except the last step; for completeness, we still present the full reduction and highlight the difference in {\color{blue}blue}:

\medskip\noindent\textbf{\underline{Reduction $\red_\mathrm{ANO\text{-}1PCA}$}:}

\medskip\noindent
On input $(pk,c)$,
\begin{enumerate}
  \item Invoke $\env$. On $(\NewSession, sid, P, P', \pw, \text{``initiator''})$ from $\env$, run $P$'s algorithm to compute $(s,T)$, except that $\red_\text{ANO-1PCA}$ uses its input $pk$ rather than $(pk,\star) \gets \Gen$. (Note that computing $(s,T)$ does not require knowing $sk$.) Send $(s,T)$ to $\adv$.
  \item On $(\NewSession, sid, P', P, \pw', \text{``respondent''})$ from $\env$ and $(s^*,T^*)$ from $\adv$, if $\neg(\pw = \pw' \land (s^*,T^*) = (s,T))$, abort. Otherwise sample $SK',\tau' \gets \{0,1\}^n$, output $SK'$ to $\env$, and send $(c,\tau')$ to $\adv$. ($c$ is part of $\red_\text{ANO-PCA1}$'s inputs.)
  \item On $(c^*,\tau^*)$ from $\adv$,
      \begin{enumerate}
        \item If $c^* = c$, do what's described in Game 1.
        \item If $c^* \neq c$, find $\adv$'s $\Htag(\pw,pk,c^*,K^*)$ query whose result is $\tau^*$. If there is none or more than one, output $SK := \bot$ to $\env$. (This matches Game 2.)
        \item \fbox{Query $\PC(sk,c^*,K^*)$}. \quad // \textcolor[gray]{0.5}{This is a valid query since we do this step only if $c^* \neq c$} \\
            If the result is 1, output $SK := \Htag(\pw,pk,c^*,K^*)$ to $\env$. \\
            If the result is 0, output $SK := \bot$ to $\env$.
      \end{enumerate}
  \item {\color{blue}When $\env$ outputs bit $b^*$, copy $\env$'s output.}
\end{enumerate}

$\red_\text{ANO-1PCA}$ uses its input $pk$ on the $P$ side and $c$ on the $P'$ side. If $b = 0$ then $c \gets \Encap(pk)$, i.e., both $P$ and $P'$ use the same $pk$, so $\red_\text{ANO-1PCA}$ simulates Game 3; if $b = 1$ then $c \gets \Encap(pk')$, i.e., $P'$ uses an independently generated $pk'$, so $\red_\text{ANO-1PCA}$ simulates Game 4. We have
\[
\Dist_\env^{3,4} \leq \Adv_{\red_\text{ANO-1PCA}}^\text{ANO-1PCA}.
\]

\emph{{\color{red}This game corresponds to game $\mathsf{G}_9$ in \cite{EPRINT:ArrBarJar25}. However, there are two differences:
\begin{enumerate}
  \item \cite{EPRINT:ArrBarJar25} reduces to ANO-PCA (where the KEM adversary can query its plaintext-checking oracle polynomially many times), rather than ANO-1PCA. Unlike in Game 3, here reducing to this stronger assumption does not provide any advantage.
  \item The assumption \cite{EPRINT:ArrBarJar25} reduces to is even further stronger in that their KEM adversary's inputs are $(pk,pk',c)$, whereas we do not give $pk'$ to the KEM adversary. Note that the reduction does not need $pk'$ at all (it only needs to send $c$ which is an encapsulation of either $pk$ or $pk'$), so the assumption in \cite{EPRINT:ArrBarJar25} is unnecessarily strong.
\end{enumerate}
}}

\paragraph{$P$ and $P'$'s passwords do not match, and $\adv$ forwards the $P$-to-$P'$ message.}
\mbox{}

\medskip\noindent\emph{Game 5:} When either $\adv$ or $P'$ makes a fresh query $H_2(\pw^*,T)$ (let the result be $t^*$)\footnote{When $P'$ makes such a query, it is in the form of $H_2(\pw',T)$ and the result is $t'$. Of course $P'$ makes this query only if $T^* = T$.}, generate $(pk^*,sk^*) \gets \Gen$, and compute
\[
r^* := s \oplus t^*, \quad R^* := T/pk^*.
\]
($(s,T)$ are the honest $P$-to-$P'$ message.) If $H_1(\pw^*,r^*)$ has already been defined and it is not $R^*$, abort. Otherwise program $H_1(\pw^*,r^*) := R^*$, and if the $H_2(\pw^*,T)$ query is made by $\adv$, also store $\left<\pw^*,pk^*,sk^*\right>$. (This in particular means that in the case where $\pw \neq \pw' \land (s^*,T^*) = (s,T)$, the generation of $pk',r',R'$ is replaced by the method described above. Note that the case where $\pw = \pw' \land (s^*,T^*) = (s,T)$ has already been addressed in Game 3 and Game 4, and $P'$ does not need to make an $H_2$ query there.)

We first upper-bound the aborting probability. For every fresh $H_2(\pw^*,T)$ query, an independent $t^* \gets \{0,1\}^{3n}$ value is sampled, which results in an independent $r^* \gets \{0,1\}^{3n}$ (regardless of the distribution of $s$). As such, there are up to $q_2+1$ independent and uniformly random $r^*$ values. The probability that one of $\adv$'s $q_1$ $H_1(\pw^*,r^*)$ query inputs happens to be any of them is upper-bounded by $q_1(q_2+1)/2^{3n}$.

Now suppose the aborting condition does not happen. The difference between Game 4 and Game 5 is that Game 4 samples $H_1(\pw^*,r^*)$ uniformly at random, whereas Game 5 defines it as $T/pk^*$. Note that $pk^*$ is not used anywhere else in the game, so the distinguishing advantage between Game 4 and Game 5 can be reduced to the PR-PK property of KEM. We only sketch the reduction $\red_\text{PR-PK}$ here: $\red_\text{PR-PK}$, on input $pk^*$ (which is either an honestly generated public key or a uniformly random value), samples $i \gets [q_2+1]$, and answers the (either $\adv$'s or an honest parties') first $i-1$ $H_2$ queries as in Game 4 and the last $q_2-i+1$ $H_2$ queries as in Game 5; for the $i$-th $H_2$ query, $\red_\text{PR-PK}$ computes
\[
r^* := s \oplus t^*, \quad R^* := T/pk^*
\]
and programs $H_1(\pw^*,r^*) := R^*$. (Note that the $\left<\pw^*,pk^*,sk^*\right>$ record is not actually used in Game 5, so $\red_\text{PR-PK}$ does not need to store it.) $\red_\text{PR-PK}$ simulates other parts of Game 4/5 honestly. Finally, $\red_\text{PR-PK}$ copies $\env$'s output bit.

A standard hybrid argument shows that (assuming the aborting condition does not happen) $\env$'s distinguishing advantage between Game 4 and Game 5 is upper-bounded by $(q_2+1)\Adv_{\red_\text{PR-PK}}^\text{PR-PK}$.\footnote{Note that here we cannot make $q_2+1$ separate reductions to the PR-PK property, and must reduce to PR-PK ``in one blow''. Furthermore, the calculation of the reduction's advantage is more complicated than the normal case where the reduction wins as long as its guess is correct. This is similar to, say, why composing a PRG polynomial times yields another PRG; see, e.g., \cite[Section~3.4.3]{book:BonSho23} for details.} Overall, we have
\[
\Dist_\env^{4,5} \leq (q_2+1)\Adv_{\red_\text{PR-PK}}^\text{PR-PK} + \frac{q_1(q_2+1)}{2^{3n}}.
\]

\emph{{\color{red}This game corresponds to game $\mathsf{G}_{10}$ in \cite{EPRINT:ArrBarJar25}, except that the aborting case was ruled out in game $\mathsf{G}_2$ in \cite{EPRINT:ArrBarJar25} (as aborting condition [A2b1]).}}

\medskip\noindent\emph{Game 6:} In the case where $\pw \neq \pw' \land (s^*,T^*) = (s,T)$, sample $SK',\tau' \gets \{0,1\}^n$.

Let $\Query_2$ be the event that $\adv$ queries \emph{either} $\Hkey(\pw',pk',c,K')$ \emph{or} $\Htag(\pw',\allowbreak pk',c,K')$. Clearly Game 5 and Game 6 are identical unless $\Query_2$ happens. We upper-bound $\Pr[\Query_2]$ via a reduction $\red_\text{OW1}$ to the OW property of KEM. This is somewhat similar to the reduction $\red_\text{OW-1PCA1}$ to the OW-1PCA property in Game 3, and we highlight the differences in {\color{blue}blue}:

\medskip\noindent\textbf{\underline{Reduction $\red_\mathrm{OW1}$}:}

\medskip\noindent
Sample $i \gets [\qkey+\qtag]$ as a guess that $\adv$'s $i$-th $\Hkey$ or $\Htag$ query triggers $\Query_2$, {\color{blue}and $j \gets [q_2+1]$ as a guess that the $j$-th $H_2$ query (by either $\adv$ or $P'$) is $H_2(\pw',T)$}. \\
On input $(pk',c)$,
\begin{enumerate}
  \item Invoke $\env$. On $(\NewSession, sid, P, P', \pw, \text{``initiator''})$ from $\env$, run $P$'s algorithm to compute $(s,T)$. Send $(s,T)$ to $\adv$.
  \item {\color{blue}Whenever there is an $H_2(\star,T)$ query (by either $\adv$ or $\red_\mathrm{OW1}$ itself on behalf of $P'$ --- see step 3 below), do what's described in Game 5. However, if this is the $j$-th query, use $\red_\mathrm{OW1}$'s input $pk'$ instead of generating $pk^*$ freshly.}
  \item On $(\NewSession, sid, P', P, \pw', \text{``respondent''})$ from $\env$ and $(s^*,T^*)$ from $\adv$, if $\neg({\color{blue}\pw \neq \pw'} \land (s^*,T^*) = (s,T))$, abort. Otherwise
      \begin{enumerate}
        {\color{blue}\item If $\adv$ has not queried $H_2(\pw',T)$, make this query on behalf of $P'$.}
        \item Sample $SK',\tau' \gets \{0,1\}^n$, output $SK'$ to $\env$, and send $(c,\tau')$ to $\adv$. ($c$ is part of $\red_\text{OW1}$'s inputs.)
      \end{enumerate}
  \item On $(c^*,\tau^*)$ from $\adv$,
      \begin{enumerate}
        \item Find $\adv$'s $\Htag(\pw,pk,c^*,K^*)$ query whose result is $\tau^*$. If there is none or more than one, output $SK := \bot$ to $\env$. (This matches Game 2.)
        \item {\color{blue}Compute $K := \Decap(sk,c)$. ($sk$ is the KEM secret key generated in step 1.)} \quad \textcolor[gray]{0.5}{// No need to query a plaintext-checking oracle since $pk$ on the $P$ side and $pk'$ on the $P'$ side are independent of each other, so the reduction can sample $sk$ on its own} \\
            If {\color{blue}$K = K^*$}, output $SK := \Hkey(\pw,pk,c^*,K)$ to $\env$. \\
            If {\color{blue}$K \neq K^*$}, output $SK := \bot$ to $\env$.
      \end{enumerate}
  \item On $\adv$'s $i$-th $\Hkey$ or $\Htag$ query, if it is in the form of $\Hkey(\pw',pk',c,(K')^*)$ or $\Htag(\pw',pk',c,(K')^*)$ for some $(K')^*$, output $(K')^*$; otherwise abort.
\end{enumerate}

Assuming $\Query_2$ happens and $\red_\text{OW1}$'s guess of index $j$ is correct, the only difference between Game 5/6 and $\red_\text{OW1}$'s simulation (until $\Query_2$ happens) is that $\red_\text{OW1}$ uses its input $pk'$ to program $H_1(\pw',r')$, and then uses its input $c \gets \Encap(pk')$ as the $P'$-to-$P$ message. Since $c$ is also computed as $\Encap(pk')$ in Game 5/6, $\red_\text{OW1}$ simulates Game 5/6 perfectly. Furthermore, if $\red_\text{OW1}$'s guess of index $i$ is correct then it wins the OW game. We have
\[
\Dist_\env^{5,6} \leq \Pr[\Query_2] \leq (q_2+1)(\qkey+\qtag)\Adv_{\red_\text{OW1}}^\text{OW}.
\]

\paragraph{Remark.}
It is very subtle why the reduction needs to guess the index $j$ and lose a factor of $q_2+1$ (and we applaud \cite{EPRINT:ArrBarJar25} for getting it right). The reason is as follows. Consider the following environment $\env$:
\begin{enumerate}
  \item Initiate $P$'s session on password $\pw$ and receive $(s,T)$.
  \item Instruct $\adv$ to make a large number of $H_2(\pw^*,T)$ queries (let the result be $t^*$), one of which is $H_2(\pw',T)$ for a special value $\pw'$.
  \item For each of the $H_2$ queries in step 2,
      \[
      \text{compute }r^* := s \oplus t^*, \quad \text{query }H_1(\pw^*,r^*).
      \]
  \item Initiate $P'$'s session on password $\pw'$ and instruct $\adv$ to forward $(s,T)$ to $P'$.
\end{enumerate}
To simulate Game 5/6 correctly, $\red_\text{OW1}$ must use its input $pk'$ to program $H_1(\pw',r') := T/pk'$ in step 3. The problem is that \emph{$\red_\mathrm{OW1}$ does not know $\pw'$ until step 4, so it has no idea which $H_1(\pw',r')$ query in step 3 needs to be programmed using $pk'$}! The only way to get around is to make a guess over all $H_2$ queries, which incurs a $q_2+1$ loss. (If we require $\env$ to always initiate $P$ and $P'$'s sessions simultaneously, this loss can be avoided since in step 3 $\red_\mathrm{OW1}$ must have already learned $\pw'$.)

\medskip\emph{{\color{red}This game corresponds to game $\mathsf{G}_{11}$ in \cite{EPRINT:ArrBarJar25}, except that they reduce to OW-PCA and avoids the $\qkey+\qtag$ loss (cf. the comment after Game 3).}}

\medskip\noindent\emph{Game 7:} Again in the case where $\pw \neq \pw' \land (s^*,T^*) = (s,T)$, recall Game 6 does the following on the $P'$ side: {\color{brown}Generate $(pk',\star) \gets \Gen$, query $t' := H_2(\pw',T)$, and compute
\[
r' := s \oplus t', \quad R' := T/pk'.
\]
($(s,T)$ are the honest $P$-to-$P'$ message.) If $H_1(\pw',r')$ has already been defined and it is not $R'$, abort. Otherwise program $H_1(\pw',r') := R'$}, \ul{compute $(c,\star) \gets \Encap(pk')$}, and sample $SK',\tau' \gets \{0,1\}^n$. Finally, output $SK'$ and send $(c,\tau')$ to $P$. (Note that $K'$ is not used anywhere since Game 6 --- in particular, $P'$ does not use $K'$ to compute $SK'$ or $\tau'$ anymore --- so we do not need to explicitly mention it.) In this game, we independently generate $(pk'',\star) \gets \Gen$ and compute $(c,\star) \gets \Encap(pk'')$ (the underlined expression).

The difference between Game 6 and Game 7 is that in Game 6 $c$ is generated by
\[
pk' = \frac{T}{H_1(\pw',s \oplus H_2(\pw',T))}
\]
which $\adv$ can compute, whereas in Game 7 $c$ is generated by an independent $pk''$. We construct a reduction $\red_\text{ANO1}$ to the ANO property of KEM. This reduction, in fact, is identical to the reduction $\red_\text{OW1}$ to the OW property in Game 6 except the last step; for completeness, we still present the full reduction and highlight the difference in {\color{blue}blue}:

\medskip\noindent\textbf{\underline{Reduction $\red_\mathrm{ANO1}$}:}

\medskip\noindent
Sample $j \gets [q_2+1]$ as a guess that the $j$-th $H_2$ query (by either $\adv$ or $P'$) is $H_2(\pw',T)$. \\
On input $(pk',c)$,
\begin{enumerate}
  \item Invoke $\env$. On $(\NewSession, sid, P, P', \pw, \text{``initiator''})$ from $\env$, run $P$'s algorithm to compute $(s,T)$. Send $(s,T)$ to $\adv$.
  \item Whenever there is an $H_2(\star,T)$ query (by either $\adv$ or $\red_\mathrm{ANO1}$ itself on behalf of $P'$ --- see step 3 below), do what's described in Game 5. However, if this is the $j$-th query, use $\red_\mathrm{ANO1}$'s input $pk'$ instead of generating $pk^*$ freshly.
  \item On $(\NewSession, sid, P', P, \pw', \text{``respondent''})$ from $\env$ and $(s^*,T^*)$ from $\adv$, if $\neg(\pw \neq \pw' \land (s^*,T^*) = (s,T))$, abort. Otherwise
      \begin{enumerate}
        \item If $\adv$ has not queried $H_2(\pw',T)$, make this query on behalf of $P'$.
        \item Sample $SK',\tau' \gets \{0,1\}^n$, output $SK'$ to $\env$, and send $(c,\tau')$ to $\adv$. ($c$ is part of $\red_\text{ANO1}$'s inputs.)
      \end{enumerate}
  \item On $(c^*,\tau^*)$ from $\adv$,
      \begin{enumerate}
        \item Find $\adv$'s $\Htag(\pw,pk,c^*,K^*)$ query whose result is $\tau^*$. If there is none or more than one, output $SK := \bot$ to $\env$. (This matches Game 2.)
        \item Compute $K := \Decap(sk,c)$. ($sk$ is the KEM secret key generated in step 1.) \quad \textcolor[gray]{0.5}{// No need to query a plaintext-checking oracle since $pk$ on the $P$ side and $pk'$ on the $P'$ side are independent of each other, so the reduction can sample $sk$ on its own} \\
            If $K = K^*$, output $SK := \Hkey(\pw,pk,c^*,K)$ to $\env$. \\
            If $K \neq K^*$, output $SK := \bot$ to $\env$.
      \end{enumerate}
  \item {\color{blue}When $\env$ outputs bit $b^*$, copy $\env$'s output.}
\end{enumerate}

Assume $\red_\text{ANO1}$'s guess of index $j$ is correct. If $b = 0$ then $c \gets \Encap(pk')$, i.e., the computation of $R'$ and $c$ use the same $pk'$, so $\red_\text{ANO-1PCA}$ simulates Game 6; if $b = 1$ then $c \gets \Encap(pk'')$, i.e., the computation of $c$ uses an independently generated $pk''$, so $\red_\text{ANO-1PCA}$ simulates Game 7. We have
\[
\Dist_\env^{6,7} \leq (q_2+1)\Adv_{\red_\text{ANO-1PCA}}^\text{ANO-1PCA}.
\]

\medskip\noindent\emph{Game 8:} Again in the case where $\pw \neq \pw' \land (s^*,T^*) = (s,T)$, we outright skip the computation of $t',r',R'$ and the programming of $H_1$ (the {\color{brown}brown} text in Game 6). (Since $pk'$ is not generated anymore, we can rename $pk''$ --- the KEM public key from which $c$ is computed --- to $pk'$.)

Note that the aborting condition, which was originally introduced in Game 5, is removed (for the specific $H_2(\pw',T)$ query by $P'$). By an analysis similar to that in Game 5, the aborting probability in Game 7 is upper-bounded by $q_1/2^{3n}$. If the aborting condition does not happen, the {\color{brown}brown} text is independent of the rest of the game, so removing it does not affect the distribution of $\env$'s output. We have
\[
\Dist_\env^{6,7} \leq \frac{q_1}{2^{3n}}.
\]

\emph{{\color{red}The two games above combined correspond to game $\mathsf{G}_{12}$ in \cite{EPRINT:ArrBarJar25}. \cite{EPRINT:ArrBarJar25} does not have the intermediate Game 7 in our proof, and outright removes the brown text from what's our Game 6. We believe it is clearer to include our Game 7, where $H_2$ is still queried and $H_1$ is still programmed: in this way it is easier to see why the reduction to ANO simulates the games correctly. (Of course, it is possible that an experienced cryptographer can see through the two game hops in one shot, but the author of this paper admits that he is not that experienced and has to work out the game hops and reductions the ``dumb'' way.) Also, \cite{EPRINT:ArrBarJar25} reduces to ANO-PCA\footnote{\cite{EPRINT:ArrBarJar25} actually mentions ANO-CPA, which appears to be a typo.}, although it also says that the plaintext-checking oracle is not needed.}}

\paragraph{Computation of the $P$-to-$P'$ message.}
\mbox{}

\medskip\noindent\emph{Game 9:} Note that up to now we have not changed $P$'s algorithm before sending $(s,T)$ (in particular, Game 5 dealt with $H_2$ queries by $\adv$ and $P'$, but not $P$). In this game we modify it as follows: Sample $s \gets \{0,1\}^{3n}, T \gets \group$. If $H_2(\pw,T)$ has already been defined, abort. Otherwise query $t := H_2(\pw,T)$, generate $(pk,sk) \gets \Gen$, and compute
\[
r := s \oplus t, \quad R := T/pk.
\]
If $H_1(\pw,r)$ has already been defined and it is not $R$, abort. Otherwise program $H_1(\pw,r) := R$.

We first upper-bound the aborting probability. Apart from the one made by $P$, there are $q_2$ $H_2$ queries in total (note that $P'$ does not query $H_2$ anymore since Game 8). Thus, the probability that one of the query inputs happens to contain $T$ is upper-bounded by $q_2/|\group|$. By an analysis similar to that in Game 5 and Game 8, the probability of the second aborting event is upper-bounded by $q_1/2^{3n}$.

Now suppose neither of the aborting events happens. The following equations hold in both Game 8 and Game 9:
\[
r = s \oplus H_2(\pw,T), \quad T = R \cdot pk.
\]
Here, $pk$ is generated by $\Gen$, and $H_2(\pw,T)$ is fixed once $T$ is fixed. The four remaining variables $r,s,R,T$ are constrained by two equations and have $4-2=2$ degrees of freedom: if we sample $r \gets \{0,1\}^{3n}, R \gets \group$ then we get Game 8, and if we sample $s \gets \{0,1\}^{3n}, T \gets \group$ then we get Game 9. Thus, Game 8 and Game 9 are identical. We have
\[
\Dist_\env^{8,9} \leq \frac{q_1}{2^{3n}} + \frac{q_2}{|\group|}.
\]

\emph{{\color{red}The game above corresponds to game $\mathsf{G}_{13}$ in \cite{EPRINT:ArrBarJar25}. We believe our argument is much more concise and cleaner.}}

\bibliographystyle{alpha}
\bibliography{bib.bib}
\end{document} 